\documentclass{article}
\usepackage{amsmath, amssymb, amsfonts}
\usepackage[utf8]{inputenc}
\usepackage[spanish]{babel}
\usepackage{hyperref}

\title{Inecuaciones}
\author{Iván Palomo y Alexandre Torrandell}
\date{23 de febrero de 2026}

\begin{document}
	\maketitle
	\tableofcontents
	\newpage
	
	\section{$|x-|x-2||<3$}
		
	\medskip
	
	Esto equivale a resolver:
	
	\begin{align*}
		|x-|x-2||<3
		&\iff -3<x-|x-2|<3\\
		&\iff 3+x > |x-2|> -3 + x. \\
	\end{align*}
	
	
	Resolvemos cada desigualdad por separado:
	
	$$
	\begin{cases}
		3+x > |x-2|,\\
		|x-2|> -3 + x
	\end{cases}
	$$ \\
	
	En la primera desigualdad distinguiremos entre otros dos subcasos:
	
		$$
	\begin{cases}
		|x-2| = x-2 \Rightarrow 3 > -2,\\
		|x-2| = -(x-2) \Rightarrow x > -\frac{1}{2}
	\end{cases}
	$$ \\
	
	En la segunda desigualdad haremos el mismo proceso:
	
		$$
	\begin{cases}
		|x-2| = x-2 \Rightarrow -2 > -3,\\
		|x-2| = -(x-2) \Rightarrow x >\frac{5}{2}
	\end{cases}
	$$ \\

	Luego,
	$$
	\boxed{ x\in(-\frac{1}{2}, \infty)}
	$$
	
	\section{$x^2 -ax \le 2$}
	
Esto equivale a resolver:
	
		$$x^2 -ax \le 2 \iff x(x-a)\le 2$$ \\
		
Los casos $x = 0$ y $x = a$ son triviales ya que $0 \le 2$ siempre, Entonces veremos que pasa con $(x-a)$.\\

Ahora queremos distinguir cuando $x<a$ y $x>a$, tomando cuando cada una es positiva o negativa. \\

Sea $x<a$ si: 
	$$
\begin{cases}
	x < a < 0 \Rightarrow \textbf{No Cumple}\\
	x < 0 < a  \Rightarrow \textbf{Cumple} \\
	0 < x < a  \Rightarrow \textbf{Cumple}
\end{cases}
$$ \\

Sea $x>a$ si: 
$$
\begin{cases}
	a < x < 0 \Rightarrow \textbf{No Cumple}\\
	a < 0 < x  \Rightarrow \textbf{Cumple}\\
	0 < a < x  \Rightarrow \textbf{Cumple}
\end{cases}
$$ \\

Luego, 
		$$
	\boxed{ x\in(a,\infty)}
	$$
	
\end{document}